% Generated draft: requires mathematical and source review.
\documentclass[11pt]{article}
\usepackage[margin=25mm]{geometry}
\usepackage{fontspec}
\setmainfont{Latin Modern Roman}
\usepackage{amsmath,amssymb,mathtools,unicode-math}
\setmathfont{Latin Modern Math}
\usepackage{enumitem,needspace,xurl,hyperref}
\setlength{\parindent}{0pt}
\setlength{\parskip}{5pt}
\setlength{\emergencystretch}{3em}
\newcommand{\N}{\mathbb{N}}
\newcommand{\Z}{\mathbb{Z}}
\newcommand{\Q}{\mathbb{Q}}
\newcommand{\R}{\mathbb{R}}
\newcommand{\C}{\mathbb{C}}
\newcommand{\F}{\mathbb{F}}
\newcommand{\A}{\mathbb{A}}
\newcommand{\PP}{\mathbb{P}}
\newcommand{\E}{\mathbb{E}}
\newcommand{\eps}{\varepsilon}
\newcommand{\dd}{\,\mathrm{d}}
\begin{document}
\section*{UnsolvedMath ID: 701400}
\textbf{Source problem code:} catalogue entry
\subsection{Definitions and mathematical statement}
A variety of groups is a class of groups defined by a collection of group identities. For a group \(G\), let \(\operatorname{Var}(G)\) denote the smallest variety containing \(G\), equivalently the class of groups satisfying every group identity valid in \(G\). A variety \(\mathcal V\) is called \emph{small} if the isomorphism classes of finitely generated groups belonging to \(\mathcal V\) form an at most countable set.

A group \(G\) satisfies the \emph{maximal condition on normal subgroups} if every ascending sequence
\[
N_1\leq N_2\leq\cdots,\qquad N_i\trianglelefteq G,
\]
is eventually constant. A group \(H\) has the \emph{Hopf property}, or is \emph{Hopfian}, if every surjective endomorphism \(H\to H\) is an automorphism.

The two questions are:
\begin{enumerate}
\item[(a)] For every finitely generated group \(G\), does smallness of \(\operatorname{Var}(G)\) imply that \(G\) satisfies the maximal condition on normal subgroups?
\item[(b)] For every variety of groups \(\mathcal V\), is the following equivalence true?
\[
\mathcal V\text{ is small}
\quad\Longleftrightarrow\quad
\text{every finitely generated group }H\in\mathcal V\text{ is Hopfian}.
\]
\end{enumerate}
The supplied registry labels these questions open; this is a supplied status, not an independently verified assessment.
\subsection{Short English statement}
P. M. Neumann asks whether countability of the finitely generated groups in a generated variety forces the generating group to satisfy the maximal condition on normal subgroups, and whether small varieties are exactly those whose finitely generated groups are all Hopfian.
\subsection{Sources}
\begin{itemize}
\item Local UnsolvedMath registry, permanent ID \(701400\), source code \texttt{catalogue entry}, title \emph{Kourovka Notebook Problem 17.98}. The supplied record attributes the problem to P. M. Neumann, proposed in 2010, and labels its status as open.
\item \emph{The Kourovka Notebook (Unsolved Problems in Group Theory)}, problem 17.98, 17th issue (2010); the supplied background also identifies the 21st issue edition (2026), \texttt{arXiv:1401.0300v48}. The supplied collection description names E. I. Khukhro and V. D. Mazurov as editors. These bibliographic details are supplied attribution and have not been independently verified.
\end{itemize}
\end{document}
